\input{preamble.tex}

\title{Лабораторная работа № 13. \\ Статическая маршрутизация в Интернете. Планирование}
\author{Сущенко Алина Николаевна \\ НПИбд-01-23}
\institute{Российский университет дружбы народов имени Патриса Лумумбы}
\date{2026}

\begin{document}

\frame{\titlepage}

\begin{frame}
\frametitle{Цель работы}
\begin{itemize}
    \item Провести подготовку по организации взаимодействия через сеть провайдера посредством статической маршрутизации локальной сети с сетью основного здания, расположенного в 42-м квартале в Москве, и сетью филиала, расположенного в г. Сочи.
\end{itemize}
\end{frame}

\begin{frame}
\frametitle{Изменение схемы сети}
    \includegraphics[width=\textwidth]{../images/01.png}
\end{frame}

\begin{frame}
\frametitle{Замена модулей на медиаконвертерах}
    \includegraphics[width=\textwidth]{../images/02.png}
\end{frame}

\begin{frame}
\frametitle{Настройка маршрутизатора}
    \includegraphics[width=\textwidth]{../images/03.png}
\end{frame}

\begin{frame}
\frametitle{Добавленный дополнительный интерфейс.}
    \includegraphics[width=\textwidth]{../images/04.png}
\end{frame}

\begin{frame}
\frametitle{Окончание настройки маршрутизатора.}
    \includegraphics[width=\textwidth]{../images/09.png}
\end{frame}

\begin{frame}
\frametitle{Схема города Москва.}
    \includegraphics[width=\textwidth]{../images/05.png}
\end{frame}

\begin{frame}
\frametitle{Размещение объектов в основном здании 42-го квартала в Москве.}
    \includegraphics[width=\textwidth]{../images/06.png}
\end{frame}

\begin{frame}
\frametitle{Размещение объектов в здании Provider.}
    \includegraphics[width=\textwidth]{../images/07.png}
\end{frame}

\begin{frame}
\frametitle{Размещение объектов в здании филиала в г. Сочи.}
    \includegraphics[width=\textwidth]{../images/08.png}
\end{frame}

\begin{frame}
\frametitle{Начальная настройка коммутатора ansusenko-q42-sw-1.}
    \includegraphics[width=\textwidth]{../images/10.png}
\end{frame}

\begin{frame}
\frametitle{Начальная настройка ansusenko-hostel-gw-1.}
    \includegraphics[width=\textwidth]{../images/11.png}
\end{frame}

\begin{frame}
\frametitle{Первоначальная настройка коммутатора ansusenko-hostel-sw-1}
    \includegraphics[width=\textwidth]{../images/12.png}
\end{frame}

\begin{frame}
\frametitle{Схема подключения подсети филиала в г. Сочи}
    \includegraphics[width=\textwidth]{../images/13.png}
\end{frame}

\begin{frame}
\frametitle{Начальная настройка sch-ansusenko-gw-1.}
    \includegraphics[width=\textwidth]{../images/14.png}
\end{frame}

\begin{frame}
\frametitle{Контрольные вопросы 1}
\textbf{1. В каких случаях следует использовать статическую маршрутизацию? Приведите примеры.}

Статическую маршрутизацию следует использовать в следующих случаях:
\begin{itemize}
\item Для небольших сетей, где не требуется высокая изменяемость маршрутов.
\item Когда нужно обеспечить предсказуемость маршрутов для трафика.
\item Для маршрутизации между филиалами, когда известны фиксированные пути.
\item Для соединения с удаленными сетями, где нет динамических изменений.
\end{itemize}

Примеры использования:
\begin{itemize}
\item Соединение главного офиса с филиалом через статически заданный маршрут.
\item Создание маршрутов к резервным интернет-провайдерам для повышения надежности.
\end{itemize}
\end{frame}

\begin{frame}
\frametitle{Контрольные вопросы 2}
\textbf{2. Укажите основные принципы статической маршрутизации между VLANs.}

Основные принципы статической маршрутизации между VLANs:
\begin{itemize}
\item Разделение сетевого трафика для разных VLAN, чтобы обеспечить безопасность и управление трафиком.
\item Настройка маршрутов вручную на маршрутизаторах или уровнях 3 коммутаторов для обмена пакетами между VLAN.
\item Использование интерфейсов под каждую VLAN и их привязка к соответствующим статическим маршрутам.
\item Проверка и корректировка маршрутов при изменении топологии сети или добавлении новых VLAN.
\end{itemize}

Пример настройки:
\begin{itemize}
\item Для VLAN 10 и VLAN 20 можно настроить маршруты:
\begin{center}
\texttt{ip route add 192.168.10.0/24 via 192.168.1.1} \\
\texttt{ip route add 192.168.20.0/24 via 192.168.1.2}
\end{center}
\end{itemize}
\end{frame}
\begin{frame}
\frametitle{Выводы}
\begin{itemize}
    \item В результате лабораторной работы мы усложнили нашу сеть, успешно провели подготовку по организации взаимодействия через сеть провайдера посредством статической маршрутизации локальной сети с сетью основного здания, расположенного в 42-м квартале в Москве, и сетью филиала, расположенного в г. Сочи.
\end{itemize}
\end{frame}

\end{document}
